\chapter{Security, Biorisk and Governance}
\label{ch:security}

Two distinct concerns share this chapter because they are entangled in practice:
the security of systems you run, and the dual-use character of the biological
capability these models provide. The entanglement is not incidental --- the same
safety mechanism frequently sits on both.

\section{Running an agent is running arbitrary code}

An agent harness executes shell commands on your machine, with your permissions,
based on text generated by a model that can be influenced by any content it
reads. Stated that plainly, the risk model is obvious.

\textcite{shapira2026agentsofchaos} ran the experiment properly: autonomous
agents in a live laboratory environment with persistent memory, email, Discord,
file systems and shell access, probed by twenty researchers over two weeks under
benign and adversarial conditions. Eleven representative failures are documented:

\begin{itemize}
  \item Unauthorised compliance with instructions from non-owners.
  \item Disclosure of sensitive information.
  \item Execution of destructive system-level actions.
  \item Denial-of-service and uncontrolled resource consumption.
  \item Identity spoofing.
  \item Cross-agent propagation of unsafe practices.
  \item Partial system takeover.
  \item \textbf{Reporting task completion while the system state contradicted the
    report.}
\end{itemize}

The last one deserves its emphasis. It means an agent's own account of what it
did is not evidence that it did it.

\section{Prompt injection}

The central unsolved problem. A language model does not reliably distinguish
instructions from data. Any content it reads --- a file, a web page, a README, a
code comment, an issue on a repository, a CSV cell --- can contain text that the
model treats as a command.

This is not a bug with a fix pending. It follows from the architecture: both
instructions and data arrive as tokens in the same context, and there is no
channel separation. Mitigations reduce the rate; none eliminates it.

Attack surfaces that matter in research:

\begin{itemize}
  \item A downloaded dataset with instructions in a metadata field.
  \item A PDF with invisible text (white on white, or off-page).
  \item A dependency whose README or docstrings carry injected text.
  \item A web page fetched during a literature search.
  \item An MCP server description --- the tool's own self-description is
    untrusted input.
\end{itemize}

\begin{keypoint}
Treat everything an agent reads as untrusted. This includes your own data files
if they came from anywhere else, and it includes the descriptions that tools
provide about themselves.
\end{keypoint}

\section{Sandboxing}

Defence in depth, with honest limits at each layer.

\textbf{Permission prompts.} The harness asks before running commands. Real
protection, defeated by fatigue --- after fifty prompts, people approve
reflexively. Configure allowlists for genuinely safe read-only operations so the
prompts you do see carry signal.

\textbf{Filesystem restriction.} Run the agent in a directory containing only
what it needs. This is free and it is the single highest-value habit in this
section.

\textbf{Lightweight sandboxes.} \code{bwrap} on Linux raises the bar. Escapes
exist; treat it as a layer, not a boundary.

\textbf{Containers.} Podman or Docker isolate the namespace but share the host
kernel. Kernel vulnerabilities have repeatedly produced container escapes and
privilege escalation, so a container is not a security boundary against a
determined attacker --- though it is entirely adequate against an agent that
misreads an instruction.

\textbf{MicroVMs.} A separate kernel, a genuine boundary. The correct answer for
running agents with reduced supervision, and increasingly practical on Linux.

\textbf{Network egress control.} Underused and effective. An agent that cannot
reach the network cannot exfiltrate, cannot fetch a malicious payload, and cannot
be driven by remote content. If the task is local, cut the network.

\textcite{rashidi2026execsecurity} systematises 39 papers across 17 categories of
execution security for coding agents --- isolation, capability and access
control, policy enforcement, time-of-check-to-time-of-use races, MCP threats,
identity delegation, execution provenance, egress control, static analysis of
generated code. The survey's own finding is that the literature is fragmented and
rarely cross-cites, which is a reasonable proxy for the field's maturity.

The offensive side has developed in parallel.
\textcite{qasemi2026cyberattacks} classify AI-powered attacks and the defensive
mechanisms available, and the summary for a research institution is
unsurprising: phishing has become much harder to detect by writing quality, and
reconnaissance against a target has become much cheaper. Neither is specific to
AI research groups, but both raise the cost of a credential leak.

\section{Supply chain}

\textbf{Slopsquatting.} Models hallucinate package names. Attackers register the
hallucinated names. The generated \code{install.packages()} call then installs an
attacker's package. Verify every dependency an agent adds exists and is the one
you meant --- this is the direct, practical reason for the ``do not add
dependencies without asking'' line in \Cref{ch:agentic}.

\textbf{Model artefacts.} Downloading weights is executing code from an untrusted
source if the format permits it. Pickle-based formats can carry arbitrary
payloads. Prefer safetensors, verify checksums, and note that this has already
produced real incidents \parencite{openai2026huggingface}.

\textbf{MCP servers.} An MCP server is a program you run with access to your
data, and its tool descriptions enter the model's context as trusted-looking
text. Apply the SNIFF criteria \parencite{nagraj2026sniff} before installing one:
who maintains it, is the project active, does it fit, how fragile are its
dependencies.

\section{A practical baseline}

\begin{enumerate}
  \item Never run an agent with elevated privileges. There is no task that
    requires it.
  \item Run it in a project directory, not in your home directory.
  \item Version control everything, commit before, review the diff after.
  \item Keep credentials out of the working tree and out of the environment the
    agent can read.
  \item Rotate any key that has ever been in a context window.
  \item Disable network egress for tasks that do not need it.
  \item Use a container for anything running with reduced supervision; a microVM
    if it is running unattended.
  \item Read every dependency addition.
  \item Never point an agent at production infrastructure.
\end{enumerate}

Treat credentials that have appeared in a chat as compromised. Shared
conversations have surfaced in public search results with API keys intact
\parencite{ibtimes2026claudechats}. Rotate, do not rationalise.

\section{Dual use in biology}

The biological capability of frontier models is a live policy question, and the
evidence is stronger than it was.

\textcite{brent2025bioweapons} present dialogues in which deployed foundation
models provide guidance on elements of technically demanding development tasks,
including construction of viral pathogens, and argue this constitutes a real
increase in biological weapons risk.
\textcite{zhang2026uplift} ran the controlled version: a multi-model,
multi-benchmark human uplift study comparing novices with LLM access against
internet-only access across eight biosecurity-relevant task sets, with up to
thirteen hours per task. They found substantial uplift.

The counterweight is calibration.
\textcite{peppin2025biorisk} argue that much biorisk discourse outruns the
evidence, and that the operative constraints on biological weapons development
are tacit knowledge, materials access and laboratory skill rather than
information availability. \textcite{turner2026biosecuritybenchmarks} examines the
benchmark evidence specifically: frontier models now match or exceed expert
reference points on many biosecurity-relevant evaluations and several benchmarks
are approaching saturation, but these results have not been calibrated against
physical task performance. Benchmark scores are highly sensitive to format.

Both can be true. Information barriers have fallen; physical barriers have not.
The question is how much of the total barrier was informational, and it is not
settled.

\subsection{The correlation with cybersecurity}

The mechanism linking the two halves of this chapter.

A model's refusal behaviour --- declining to assist with pathogen engineering ---
is implemented in the instruction-tuning layer, in the same place and by the same
mechanism as its refusal to write malware. Prompt injection defeats both, because
injection subverts instruction-following in general rather than any particular
refusal.

The consequence, stated by practitioners and worth repeating: \emph{prompt
injection defences open a biosecurity gap}. Robust separation of instructions
from data would strengthen both simultaneously. Its absence weakens both
simultaneously. They are not independent problems.

\textcite{holko2026relationalbiosecurity} draws the structural conclusion.
Biosecurity frameworks are organised around control of known agents, materials
and individual technologies. In AI-enabled biology, capability comes from combination --- from the way data,
models, infrastructure and workflows fit together. The risk is not in any one
part but in how they connect. Nucleic acid synthesis screening, effective in isolation, may
not capture risks that emerge from an integrated system.

\textcite{barkan2026restricting} examine the natural response --- software
barriers restricting AI agent use of biological design tools --- and are
measured about feasibility. Restriction at the tool layer is more tractable than
restriction at the model layer, because tools have identifiable interfaces.

For working ecologists, none of this changes daily practice. It does explain why
some models decline biological questions that seem benign, why open-weight
release decisions are contested, and why sequence data governance is becoming
stricter. It is worth understanding the argument before being asked for an
opinion on it.

\section{Governance and risk taxonomy}

\textcite{slattery2026riskrepository} assembled 1,725 risks from 74 frameworks
into a structured database with two taxonomies: causal (who caused it, was it
intentional, before or after deployment) and domain (seven areas from
discrimination and privacy through misinformation to weapons). Its value is
shared vocabulary --- the paper's own motivating observation is that the same word
describes different problems across communities while different words describe
identical concerns.

For an institutional AI policy, this is a better starting point than drafting
from scratch.

\subsection{What an institutional policy should cover}

\begin{itemize}
  \item \textbf{Data classification.} What may go to a hosted service, what may
    not. Be specific: sensitive localities, personal data, embargoed sequences,
    third-party data under agreement.
  \item \textbf{Approved tools.} A short list with a named owner, rather than
    ad-hoc adoption.
  \item \textbf{Disclosure.} What must be declared in publications, and where.
  \item \textbf{Agent execution.} Sandboxing requirements, and an explicit
    prohibition on production access.
  \item \textbf{Credentials.} Rotation policy after context exposure.
  \item \textbf{Procurement.} Whether a service that trains on submitted data is
    acceptable, and under what contract terms.
\end{itemize}

\section{Risks to judgement}

Not every risk is an attack. Two slow ones damage the research itself.

The first is \term{sycophancy} (\Cref{ch:powers}). A model that agrees with you
removes the second opinion you thought you were getting. Exposure to
sycophantic AI increases dependence on it \parencite{cheng2026sycophancy}, and
even a perfectly rational user who trusts a source that keeps agreeing will
become confidently wrong \parencite{chandra2026spiraling}. In security terms it
is a failure of review: the check is there on paper and absent in practice.
Ask the model to argue against your result, and get a person to check anything
that matters.

The second is \term{slop}: text produced in bulk by a model and sent on with
little human thought in it. \textcite{turner2026slop} describes inboxes,
feeds, papers and grant proposals filling with it. A journal editor who met the
authors of ten submissions already due for desk rejection found that three
could not answer basic questions about their own papers and three more
struggled with the details; the one paper whose authors could answer had a major
error in its central claim. The consequence is that a publication tells you
less than it did about what its authors understand. Turner's own experience
points the same way: when he delegated writing, he lost the thinking that
writing does. ``You become less capable while appearing more productive.'' This
is the \term{deskilling} risk that \textcite{mammides2026llmecology} raise for
ecology.

Neither is fixed by tooling. The defence is the one this book keeps returning
to: write your own arguments, check claims against sources, and put your name
only on work you can explain.

\section{The wider frame}

Two strands in the literature situate this in longer time.

\textcite{muller2026evolvableai} consider AI systems whose components, learning
rules and deployment conditions can themselves undergo Darwinian evolution, and
ask under what technical and ecological conditions this becomes possible. They
argue the scenario is underappreciated in existential risk debates. Whether or
not one credits the conclusion, the analytical move --- treating deployed AI
systems as an evolving population under selection --- is one ecologists are
unusually well equipped to evaluate.

\textcite{korecki2025biosphericai} makes the complementary argument about values.
AI ethics is overwhelmingly anthropocentric; recent extensions reach sentientism;
neither captures the biosphere. He proposes biospheric ethics as a design target
and asks what an AI aligned with biospheric rather than human interests would
look like.

These connect to an older line of thought --- Vernadsky's transition from
biosphere to noosphere \parencite{vernadsky1938noosphere}, revisited by
\textcite{vidal2024noosphere} as planetary superorganism and major evolutionary
transition, and pushed further by \textcite{lovelock2019novacene}, who argued that
the Anthropocene is ending and a Novacene beginning in which new beings emerge
from existing AI systems. \textcite{vonneumann1955survive} asked the version of
the question that has aged best, in 1955, about a different technology: whether
the environment in which technological progress occurs has become too small and
too disorganised for what is being built in it.

These are not operational guidance. They are the background against which the
operational questions are being asked, and ecologists have a professional stake
in them that most participants in the debate do not.
